\documentclass[12pt,tikz,border=5pt]{article}
\usepackage[margin=1in]{geometry}
\usepackage[spanish]{babel}
\addto\captionsspanish{\renewcommand{\tablename}{Tabla}}
\decimalpoint
\usepackage{tikz}
\usepackage{subcaption}
\usetikzlibrary{babel} 
\usepackage{setspace}
\usepackage{booktabs}
\usepackage{natbib}   
\usepackage{hyperref}
\usepackage{enumitem}
\usepackage{multicol}
\usepackage{amsmath}
\usepackage{amssymb} 
\usepackage{amsfonts}
\usepackage{amsthm}
\newtheorem{ejercicio}{Ejercicio}[section]

\doublespacing

\title{Cálculo Diferencial e Integral \\ Licenciatura en Ingeniería en Computación}
\author{M. en C. Edwin Pérez \\ Dr. Julian Bravo}
\date{\today}

\begin{document}

\maketitle

\section{Maneras de Representar una Función}
Intenta resolver los ejercicios propuestos antes de la clase del 10 de septiembre. Lleva tus dudas para discutirlos en clase. Los ejercicios marcados como tarea se entregarán el 17 de septiembre en la plataforma classroom y escritos preferentemente en latex. \\

Ejercicios tomados del libro de Cálculo de una Variable Trascendentes Tempranas 4ta Edición\citep{stewart2001calculo} (la sección del libro donde vienen los ejercicios para la tarea se puede encontrar en el classroom, dentro de la sección Tarea\,\#1, está nombrada como ``Ejercicios\_Representacion\_Funciones.pdf''.

\begin{ejercicio}
  Resuelve lo siguiente y exprese claramente su procedimiento.
  \begin{enumerate}[label=(\alph*)]
    \item Halle la pendiente de la línea tangente a la parábola $y = 4x - x^2$ en el punto $(1, 3)$.
      \begin{enumerate}
        \item usando la definición de derivada con incrementos, $$\displaystyle f'(a) = \lim_{\Delta t \to 0} \frac{f(a+\Delta t)-f(a)}{\Delta t}$$
        \item usando la forma equivalente, $$\displaystyle f'(a) = \lim_{x \to a} \frac{f(x)-f(a)}{x-a}$$
      \end{enumerate}
    \item Encuentre una ecuación de la recta tangente del inciso (a).
    \item Dibuje la parábola y la tangente.
  \end{enumerate}
\end{ejercicio}

\begin{ejercicio}
  Resuelve lo siguiente y exprese claramente su procedimiento.
  \begin{enumerate}[label=(\alph*)]
    \item Encuentre la pendiente de la tangente a la curva $y = x - x^3$ en el punto $(1, 0)$.
      \begin{enumerate}
        \item usando la definición de derivada con incrementos, $$\displaystyle f'(a) = \lim_{\Delta t \to 0} \frac{f(a+\Delta t)-f(a)}{\Delta t}$$
        \item usando la forma equivalente, $$\displaystyle f'(a) = \lim_{x \to a} \frac{f(x)-f(a)}{x-a}$$
      \end{enumerate}
    \item Halle una ecuación de la recta tangente del inciso (a).
    \item Dibuje la curva y la tangente.
  \end{enumerate}
\end{ejercicio}

\begin{ejercicio}
  En los problemas siguientes, damos la función de posición $x = f(t)$ de una partícula en movimiento en una recta horizonta. Determine su posición $x$ cuando su velocidad $v$ es cero.
  \begin{enumerate}[label=(\alph*)]
    \item $x = 100 - 16t^2$
    \item $x = -16t^2 + 80t - 1$
    \item $100- 20t - 5t^2$
  \end{enumerate}
\end{ejercicio}

\begin{ejercicio}
  En los problemas siguientes, damos la altura $y(t)$ (en metros en el instante $t$ segundos) de una pelota lanzada verticalmente hacia arriba. Determine la altura máxima que alcanza la pelota.
  \begin{enumerate}[label=(\alph*)]
    \item $y = -16t^2 + 160t$
    \item $y = -16t^2 + 128t + 25$
    \item $y = -16t^2 + 96t + 50$
  \end{enumerate}
\end{ejercicio}

\begin{ejercicio}
  La temperatura Celsius $C$ está dada en términos de la temperatura Fahrenheit $F$ por $C = \dfrac{5}{9}(F - 32)$. Determine la razón de cambio de $C$ con respecto de $F$ y la razón de cambio de $F$ con respecto de $C$.
\end{ejercicio}

\begin{ejercicio}
  Determine la razón de cambio del área de $A$ de un círculo con respecto de su circunferencia $C$.
\end{ejercicio}

\begin{ejercicio}
  Determine la razón de cambio del área de $A$ de un círculo con respecto de su circunferencia $C$.
\end{ejercicio}

\begin{ejercicio}
  Un globo esférico con un radio inicial $r = 5$ cm comineza a deinflarse en el instante $t=0$ y su radio $t$ segundos más tarde es $r =  \dfrac{60 - t}{12}$ cm. ¿A qué razón (en $cm^3/s$) sale el aire del globo cuando $t = 30$?
\end{ejercicio}

\begin{ejercicio}
  El volumen $V$ (en litros) de 3 gramos de $CO_2$ a 27$^oC$ está dado en términos de su presion $p$ (en atmósferas) por la fórmula $V(p) = \dfrac{1.68}{p}$. ¿Cuál es la razón de cambio de $V$ con respecto de $p$ cuando $p = 2$ (atmósferas)?
\end{ejercicio}

\begin{ejercicio}
  Al derretirse una bola de nieve con radio inicial de 12 cm, su radio decrece a una razón constante. Comienza a derretirse cuando $t = 0$ h (horas) y tarda 12 horas en desaparecer.
  \begin{enumerate}[label=(\alph*)]
    \item ¿Cuál es la razón de cambio del volumen cuando $t = 6$
    \item ¿Cuál es la razón de cambio promedio del volumen de $t = 3$ a $t = 9$.
  \end{enumerate}
\end{ejercicio}

\begin{ejercicio}
  En los siguientes ejercicios, encuentre la pendiente de la recta tangeten a la gráfica de la función en el punto dado.
  \begin{multicols}{3}
    \begin{enumerate}[label=(\alph*)]
      \item $f(x) = 3 - 5x$, $(-1, 8)$
      \item $g(x) = x^2 - 9$, $(2, -5)$
      \item $f(t) = 3t - t^2$, $(0, 0)$
      \item $g(x) = \dfrac{3}{2}x + 1$, $(-2, -2)$
      \item $f(x) = 5 - x^2$, $(3, -4)$
      \item $h(t) = t^2 + 4t$, $(1, 5)$
    \end{enumerate}
  \end{multicols}
\end{ejercicio}

\begin{ejercicio}
  En los siguientes ejercicios, encuentre la pendiente de la recta tangeten a la gráfica de la función en el punto dado.
  \begin{multicols}{3}
    \begin{enumerate}[label=(\alph*)]
      \item $f(x) = 3 - 5x$, $(-1, 8)$
      \item $g(x) = x^2 - 9$, $(2, -5)$
      \item $f(t) = 3t - t^2$, $(0, 0)$
      \item $g(x) = \dfrac{3}{2}x + 1$, $(-2, -2)$
      \item $f(x) = 5 - x^2$, $(3, -4)$
      \item $h(t) = t^2 + 4t$, $(1, 5)$
    \end{enumerate}
  \end{multicols}
\end{ejercicio}

\bibliographystyle{abbrv}
\bibliography{cite}

\end{document}
